\chapter{Lezione 22 --- Controllo decentralizzato e disaccoppiamento nei sistemi MIMO}

\section{Richiamo: margine di fase e margine di guadagno}

Nel diagramma di Nyquist, la stabilità del sistema in ciclo chiuso si studia rispetto al punto critico
\[
-1.
\]
Le due quantità classiche che misurano la ``distanza'' dal punto critico sono:

\begin{itemize}
    \item \textbf{margine di fase} \(M_\varphi\),
    \item \textbf{margine di guadagno} \(M_G\).
\end{itemize}

\begin{figure}[h!]
\centering
\begin{tikzpicture}[scale=1.1,>=Latex]
    % assi
    \draw[->] (-2.3,0)--(2.8,0) node[right] {};
    \draw[->] (0,-2.4)--(0,2.3) node[above] {};

    % punto critico
    \fill (-1,0) circle (1.3pt);
    \node[below] at (-1,0) {$-1$};

    % cerchi qualitativi
    \draw[thick] (0,0) circle (1.35);
    \draw[thick] (0,0) circle (0.75);

    % curve Nyquist qualitative
    \draw[very thick,red]
        (2.4,0) .. controls (1.5,-1.8) and (0.3,-2.1) .. (-1.1,-1.2)
        .. controls (-1.8,-0.8) and (-1.2,-0.1) .. (-0.35,0.15)
        .. controls (0.4,0.25) and (0.8,-0.05) .. (0.55,-0.45)
        .. controls (0.3,-0.8) and (-0.2,-0.55) .. (-0.05,-0.1)
        .. controls (0.1,0.15) and (0.35,0.12) .. (0.48,0.0);

    \draw[very thick]
        (2.1,0) .. controls (1.3,-1.5) and (0.2,-1.8) .. (-0.95,-1.05)
        .. controls (-1.4,-0.65) and (-1.0,-0.1) .. (-0.3,0.05)
        .. controls (0.25,0.15) and (0.55,-0.05) .. (0.38,-0.33)
        .. controls (0.2,-0.58) and (-0.08,-0.38) .. (0.0,-0.08)
        .. controls (0.08,0.08) and (0.22,0.05) .. (0.3,0.0);

    % frecce direzione
    \draw[->,thick] (-1.15,-1.15)--(-1.32,-1.0);
    \draw[->,thick] (-0.55,-0.55)--(-0.7,-0.4);
    \draw[->,thick,red] (-1.3,-1.25)--(-1.5,-1.1);
    \draw[->,thick,red] (-0.68,-0.72)--(-0.82,-0.56);

    % margine di guadagno
    \draw[dashed] (-1,0)--(-1,1.5);
    \draw[dashed] (-1.45,0)--(-1.45,1.5);
    \draw[<->,blue] (-1.45,1.38)--(-1,1.38);
    \node[blue,above] at (-1.225,1.42) {\small \(1/M_G\)};

    % margine di fase
    \draw[blue,thick] (-1,0)--(-1.0,-0.85);
    \draw[blue,thick] (-1,0)--(-0.42,-0.55);
    \draw[->,blue,thick] (-0.86,-0.35) arc[start angle=-100,end angle=-52,radius=0.52];
    \node[blue] at (-0.92,-0.72) {\small \(M_\varphi\)};
\end{tikzpicture}
\caption{Interpretazione qualitativa di margine di fase \(M_\varphi\) e margine di guadagno \(M_G\) sul diagramma di Nyquist.}
\end{figure}

\section{Sistema MIMO \(2\times 2\)}

Consideriamo un impianto MIMO \(2\times 2\) con matrice di trasferimento
\[
G(s)=
\begin{bmatrix}
G_{11}(s) & G_{12}(s)\\
G_{21}(s) & G_{22}(s)
\end{bmatrix}.
\]

I singoli elementi hanno il significato seguente:

\begin{itemize}
    \item \(G_{11}(s)\): funzione di trasferimento tra ingresso \(1\) e uscita \(1\),
    \item \(G_{22}(s)\): funzione di trasferimento tra ingresso \(2\) e uscita \(2\),
    \item \(G_{21}(s)\): funzione di trasferimento tra ingresso \(1\) e uscita \(2\),
    \item \(G_{12}(s)\): funzione di trasferimento tra ingresso \(2\) e uscita \(1\).
\end{itemize}

I termini fuori diagonale \(G_{12}\) e \(G_{21}\) rappresentano le \textbf{interferenze tra canali}.

\section{Controllo decentralizzato}

Nel controllo decentralizzato si usano due regolatori separati, uno per ciascun canale, ignorando nella sintesi iniziale i termini di accoppiamento.

\[
R(s)=
\begin{bmatrix}
R_1(s) & 0\\
0 & R_2(s)
\end{bmatrix}.
\]

La struttura a blocchi è la seguente.

\begin{figure}[h!]
\centering
\begin{tikzpicture}[scale=0.95,>=Latex]
    % r1 path
    \node at (-0.2,2.8) {$r_1$};
    \draw (-0.1,2.8) circle (0.08);
    \draw[->] (0,2.8)--(0.8,2.8);
    \node at (1.0,2.8) {$\sum$};
    \draw[->] (1.2,2.8)--(2.1,2.8);
    \draw (2.1,2.45) rectangle (3.2,3.15);
    \node at (2.65,2.8) {$R_1$};

    \draw[->] (3.2,2.8)--(4.2,2.8);
    \node at (4.4,2.8) {$\sum$};
    \draw[->] (4.6,2.8)--(5.5,2.8);
    \node[above] at (4.95,2.9) {$u_1$};
    \draw (5.5,2.45) rectangle (6.8,3.15);
    \node at (6.15,2.8) {$G_{11}(s)$};
    \draw[->] (6.8,2.8)--(8.0,2.8);
    \node[right] at (8.0,2.8) {$y_1$};

    % feedback y1
    \draw (7.6,2.8)--(7.6,4.1)--(0.8,4.1)--(0.8,2.92);
    \draw[->] (0.8,2.92)--(0.8,2.88);

    % r2 path
    \node at (-0.2,0.8) {$r_2$};
    \draw (-0.1,0.8) circle (0.08);
    \draw[->] (0,0.8)--(0.8,0.8);
    \node at (1.0,0.8) {$\sum$};
    \draw[->] (1.2,0.8)--(2.1,0.8);
    \draw (2.1,0.45) rectangle (3.2,1.15);
    \node at (2.65,0.8) {$R_2$};

    \draw[->] (3.2,0.8)--(4.2,0.8);
    \node at (4.4,0.8) {$\sum$};
    \draw[->] (4.6,0.8)--(5.5,0.8);
    \node[above] at (4.95,0.9) {$u_2$};
    \draw (5.5,0.45) rectangle (6.8,1.15);
    \node at (6.15,0.8) {$G_{22}(s)$};
    \draw[->] (6.8,0.8)--(8.0,0.8);
    \node[right] at (8.0,0.8) {$y_2$};

    % feedback y2
    \draw (7.6,0.8)--(7.6,-0.5)--(0.8,-0.5)--(0.8,0.68);
    \draw[->] (0.8,0.68)--(0.8,0.72);

    % interaction blocks
    \draw (4.3,1.45) rectangle (5.2,2.05);
    \node at (4.75,1.75) {$G_{12}$};
    \draw[->] (4.75,1.15)--(4.75,1.45);
    \draw[->] (4.75,2.05)--(4.75,2.6);

    \draw (3.5,1.45) rectangle (4.4,2.05);
    \node at (3.95,1.75) {$G_{21}$};
    \draw[->] (3.95,2.45)--(3.95,2.05);
    \draw[->] (3.95,1.45)--(3.95,1.0);

    \node at (8.6,2.1) {\small controllo decentralizzato};
\end{tikzpicture}
\caption{Struttura di controllo decentralizzato per un sistema \(2\times 2\): i canali principali sono regolati separatamente, ma restano le interferenze incrociate \(G_{12}\) e \(G_{21}\).}
\end{figure}

Il controllo decentralizzato, da solo, \textbf{non elimina} le interferenze tra canali.

\section{Inserimento del disaccoppiatore}

Per ridurre o annullare le interazioni, si introduce un sistema dinamico detto \textbf{disaccoppiatore}
\[
D(s),
\]
posto tra il regolatore e l'impianto.

Lo schema complessivo diventa

\[
R(s)\;\longrightarrow\; D(s)\;\longrightarrow\; G(s).
\]

Definiamo
\[
M(s)=G(s)D(s).
\]

L'obiettivo è scegliere \(D(s)\) in modo che \(M(s)\) risulti \textbf{diagonale}, o almeno quasi diagonale. Se questo accade, il sistema diventa disaccoppiato e il controllo decentralizzato diventa naturale.

\begin{figure}[h!]
\centering
\begin{tikzpicture}[scale=1,>=Latex]
    \draw (0,0) circle (0.08);
    \draw[->] (0.08,0)--(1.0,0);
    \node at (0.5,0.25) {$e$};
    \draw (1.0,-0.35) rectangle (2.0,0.35);
    \node at (1.5,0) {$R(s)$};
    \draw[->] (2.0,0)--(3.0,0);
    \draw (3.0,-0.35) rectangle (4.0,0.35);
    \node at (3.5,0) {$D(s)$};
    \draw[->] (4.0,0)--(5.0,0);
    \draw (5.0,-0.35) rectangle (6.0,0.35);
    \node at (5.5,0) {$G(s)$};
    \draw[->] (6.0,0)--(7.0,0);
    \node[right] at (7.0,0) {$y$};

    \draw (6.6,0)--(6.6,-1.0)--(0.5,-1.0)--(0.5,-0.08);
    \draw[->] (0.5,-0.08)--(0.5,-0.02);
    \node at (6.9,0.25) {};
\end{tikzpicture}
\caption{Struttura di controllo con disaccoppiatore dinamico \(D(s)\).}
\end{figure}

\section{Richiamo sul controllo classico}

Nel caso SISO classico, con
\[
e=y^\star-y,\qquad u=Re,\qquad y=Gu,
\]
si ha
\[
y=GR(y^\star-y)=GRy^\star-GRy.
\]
Da cui
\[
(I+GR)y=GR\,y^\star
\]
e quindi
\[
y=(I+GR)^{-1}GR\,y^\star.
\]

La matrice/funzione
\[
H=(I+GR)^{-1}GR
\]
è la funzione di trasferimento in ciclo chiuso.

Nel caso MIMO con disaccoppiatore, si desidera che il blocco equivalente
\[
M(s)=G(s)D(s)
\]
sia diagonale.

\section{Disaccoppiamento ideale nel caso \(2\times 2\)}

Supponiamo di volere
\[
M(s)=
\begin{bmatrix}
M_{11}(s) & 0\\
0 & M_{22}(s)
\end{bmatrix}.
\]
Se \(G(s)\) è invertibile, allora
\[
D(s)=G^{-1}(s)M(s).
\]

Per una matrice \(2\times 2\),
\[
G^{-1}(s)=\frac{1}{\Delta(s)}
\begin{bmatrix}
G_{22}(s) & -G_{12}(s)\\
-G_{21}(s) & G_{11}(s)
\end{bmatrix},
\]
dove
\[
\Delta(s)=G_{11}(s)G_{22}(s)-G_{12}(s)G_{21}(s).
\]

Quindi
\[
G^{-1}(s)M(s)
=\frac{1}{\Delta(s)}
\begin{bmatrix}
G_{22}M_{11} & -G_{12}M_{22}\\
-G_{21}M_{11} & G_{11}M_{22}
\end{bmatrix}.
\]

Una scelta naturale è mantenere sulle diagonali le dinamiche principali originarie:
\[
M_{11}(s)=G_{11}(s),\qquad M_{22}(s)=G_{22}(s).
\]
Con questa scelta si ottiene il \textbf{disaccoppiatore ideale}
\[
\boxed{
D(s)=
\begin{bmatrix}
\dfrac{G_{22}G_{11}}{G_{11}G_{22}-G_{12}G_{21}}
&
-\dfrac{G_{12}G_{22}}{G_{11}G_{22}-G_{12}G_{21}}
\\[1.2em]
-\dfrac{G_{21}G_{11}}{G_{11}G_{22}-G_{12}G_{21}}
&
\dfrac{G_{11}G_{22}}{G_{11}G_{22}-G_{12}G_{21}}
\end{bmatrix}.
}
\]

Ed effettivamente
\[
G(s)D(s)=
\begin{bmatrix}
G_{11}(s) & 0\\
0 & G_{22}(s)
\end{bmatrix}.
\]

\section{Interpretazione del disaccoppiamento ideale}

Formalmente il risultato sembra perfetto: il sistema risultante è diagonale. Tuttavia questo disaccoppiamento avviene tramite \textbf{cancellazioni dinamiche}.

Infatti
\[
M(s)=G(s)D(s)=G(s)G^{-1}(s)\,\operatorname{diag}(G(s))
=\operatorname{diag}(G(s)).
\]

Quindi stiamo cancellando, attraverso \(D(s)\), le parti non diagonali e parte delle dinamiche del sistema.

Dal punto di vista fisico questo è delicato, perché le cancellazioni possono eliminare modi interni che risultano poi non raggiungibili o non osservabili nello schema equivalente.

\section{Quando si possono cancellare poli e zeri?}

Nel progetto del regolatore si possono cancellare soltanto poli e zeri \textbf{a parte reale strettamente negativa}, cioè associati a \textbf{fase minima}.

Quindi, poiché il disaccoppiamento ideale lavora proprio per cancellazioni, è necessario richiedere che:

\begin{itemize}
    \item le dinamiche coinvolte siano a fase minima;
    \item non ci siano ritardi di tempo significativi;
    \item le cancellazioni siano fisicamente robuste.
\end{itemize}

In sintesi:

\[
\boxed{\text{il disaccoppiatore ideale richiede condizioni di fase minima}}
\]

ed è inoltre molto sensibile a:

\begin{itemize}
    \item incertezze parametriche,
    \item errori di modello,
    \item ritardi di tempo.
\end{itemize}

\section{Caso quasi diagonale}

Se i termini fuori diagonale sono piccoli, non è necessario disaccoppiare in modo esatto.

In particolare, se
\[
G_{11}(s)G_{22}(s)\gg G_{12}(s)G_{21}(s),
\]
allora nel determinante
\[
\Delta(s)=G_{11}(s)G_{22}(s)-G_{12}(s)G_{21}(s)
\]
il termine dominante è \(G_{11}G_{22}\), mentre il prodotto \(G_{12}G_{21}\) può essere trascurato.

In questo caso il sistema è \textbf{quasi diagonale}.

\section{Disaccoppiatore semplificato}

Trascurando il prodotto \(G_{12}G_{21}\) nel denominatore, dal disaccoppiatore ideale si ottiene il \textbf{disaccoppiatore semplificato}
\[
\boxed{
D(s)\approx
\begin{bmatrix}
1 & -\dfrac{G_{12}(s)}{G_{11}(s)}\\[0.8em]
-\dfrac{G_{21}(s)}{G_{22}(s)} & 1
\end{bmatrix}.
}
\]

Questo non disaccoppia perfettamente il sistema, ma riduce sensibilmente l'interazione cross-channel.

\subsection{Caso triangolare}

Se il sistema è triangolare, ad esempio
\[
G(s)=
\begin{bmatrix}
G_{11}(s) & 0\\
G_{21}(s) & G_{22}(s)
\end{bmatrix},
\]
si può scegliere
\[
D(s)=
\begin{bmatrix}
1 & 0\\[0.4em]
-\dfrac{G_{21}(s)}{G_{22}(s)} & 1
\end{bmatrix}.
\]

In questo caso la cancellazione coinvolge solo il termine \(G_{21}\) rispetto a \(G_{22}\), e la struttura è più favorevole del caso generale.

\section{Schema a blocchi con disaccoppiatore semplificato}

Il disaccoppiatore semplificato introduce due rami di compensazione incrociati che cercano di annullare, almeno in prima approssimazione, le interferenze.

\begin{figure}[h!]
\centering
\begin{tikzpicture}[scale=0.9,>=Latex]
    % upper channel
    \node at (-0.3,3.2) {$r_1$};
    \draw (-0.15,3.2) circle (0.08);
    \draw[->] (-0.07,3.2)--(0.7,3.2);
    \node at (0.9,3.2) {$\sum$};
    \draw[->] (1.1,3.2)--(2.0,3.2);
    \draw (2.0,2.85) rectangle (3.1,3.55);
    \node at (2.55,3.2) {$R_1(s)$};
    \draw[->] (3.1,3.2)--(4.2,3.2);

    \node at (4.4,3.2) {$\sum$};
    \draw[->] (4.6,3.2)--(5.5,3.2);
    \draw (5.5,2.85) rectangle (6.8,3.55);
    \node at (6.15,3.2) {$G_{11}(s)$};
    \draw[->] (6.8,3.2)--(8.0,3.2);
    \node[right] at (8.0,3.2) {$y_1$};

    % lower channel
    \node at (-0.3,0.8) {$r_2$};
    \draw (-0.15,0.8) circle (0.08);
    \draw[->] (-0.07,0.8)--(0.7,0.8);
    \node at (0.9,0.8) {$\sum$};
    \draw[->] (1.1,0.8)--(2.0,0.8);
    \draw (2.0,0.45) rectangle (3.1,1.15);
    \node at (2.55,0.8) {$R_2(s)$};
    \draw[->] (3.1,0.8)--(4.2,0.8);

    \node at (4.4,0.8) {$\sum$};
    \draw[->] (4.6,0.8)--(5.5,0.8);
    \draw (5.5,0.45) rectangle (6.8,1.15);
    \node at (6.15,0.8) {$G_{22}(s)$};
    \draw[->] (6.8,0.8)--(8.0,0.8);
    \node[right] at (8.0,0.8) {$y_2$};

    % feedbacks
    \draw (7.6,3.2)--(7.6,4.3)--(0.7,4.3)--(0.7,3.28);
    \draw[->] (0.7,3.28)--(0.7,3.24);

    \draw (7.6,0.8)--(7.6,-0.3)--(0.7,-0.3)--(0.7,0.72);
    \draw[->] (0.7,0.72)--(0.7,0.76);

    % decoupler blocks
    \draw (3.55,1.75) rectangle (4.8,2.35);
    \node at (4.175,2.05) {$-\dfrac{G_{12}}{G_{11}}$};
    \draw[->] (4.175,1.15)--(4.175,1.75);
    \draw[->] (4.175,2.35)--(4.175,3.0);

    \draw (3.0,1.75) rectangle (4.2,2.35);
    \node at (3.6,2.05) {$-\dfrac{G_{21}}{G_{22}}$};
    \draw[->] (3.6,2.85)--(3.6,2.35);
    \draw[->] (3.6,1.75)--(3.6,1.0);

    % cross couplings remaining
    \draw (5.0,1.75) rectangle (5.8,2.35);
    \node at (5.4,2.05) {$G_{21}$};
    \draw[->] (5.4,2.85)--(5.4,2.35);
    \draw[->] (5.4,1.75)--(5.4,1.0);

    \draw (5.9,1.75) rectangle (6.7,2.35);
    \node at (6.3,2.05) {$G_{12}$};
    \draw[->] (6.3,1.15)--(6.3,1.75);
    \draw[->] (6.3,2.35)--(6.3,3.0);

    % ellipse around decoupler
    \draw[green!60!black,thick] (3.95,2.0) ellipse (1.35 and 1.5);
    \node[green!60!black] at (2.9,2.95) {$D(s)$};
\end{tikzpicture}
\caption{Schema qualitativo con disaccoppiatore semplificato: l'interazione non viene eliminata del tutto, ma viene ridotta.}
\end{figure}

\section{Generalizzazione qualitativa al caso \(n\times n\)}

Nel caso \(n\times n\), l'idea del disaccoppiatore semplificato è costruire una matrice del tipo
\[
D(s)\approx
\begin{bmatrix}
1 & -\dfrac{G_{12}}{G_{11}} & \cdots & -\dfrac{G_{1n}}{G_{11}}\\[0.8em]
-\dfrac{G_{21}}{G_{22}} & 1 & \cdots & -\dfrac{G_{2n}}{G_{22}}\\
\vdots & \vdots & \ddots & \vdots\\
-\dfrac{G_{n1}}{G_{nn}} & -\dfrac{G_{n2}}{G_{nn}} & \cdots & 1
\end{bmatrix}.
\]

Questa espressione è utile come linea guida, ma non sempre può essere realizzata in pratica: bisogna verificare che non compaiano cancellazioni non ammissibili, ritardi, o termini instabili.

\section{Osservazioni finali}

\begin{itemize}
    \item Il \textbf{controllo decentralizzato} usa un regolatore per ciascun canale, ma lascia le interazioni.
    \item Il \textbf{disaccoppiatore} viene inserito per trasformare l'impianto in uno quasi diagonale o diagonale.
    \item Il \textbf{disaccoppiatore ideale} si ottiene imponendo
    \[
    M(s)=G(s)D(s)
    \]
    diagonale e usando
    \[
    D(s)=G^{-1}(s)M(s).
    \]
    \item Questa soluzione è però basata su \textbf{cancellazioni dinamiche}, quindi richiede condizioni di fase minima ed è molto sensibile a ritardi e incertezze.
    \item Se i canali sono già debolmente accoppiati, cioè se
    \[
    G_{11}G_{22}\gg G_{12}G_{21},
    \]
    conviene usare il \textbf{disaccoppiatore semplificato}, più robusto e più facile da realizzare.
\end{itemize}